% Part: first-order-logic
% Chapter: axiomatic-deduction
% Section: identity
%READERNOTE{OLAXD-009: गोठवलेल्या इंग्रजी स्रोताच्या अध्याय-शीर्ष टिप्पणीमध्ये जुने axiomatic-proofs नाव आहे; मार्ग आणि olfileid यांच्याशी सुसंगत असे axiomatic-deduction नाव येथे वापरले आहे.}

\documentclass[../../../include/open-logic-section]{subfiles}

\begin{document}

\olfileid{fol}{axd}{ide}

\olsection{\usetoken{P}{derivation} सह \usetoken{S}{identity}}

!!{derivation}s मध्ये $\eq$ समाविष्ट करण्यासाठी आपण फक्त नवीन
स्वयंसिद्धक-आकार जोडतो. !!{derivation} आणि $\Proves$ यांच्या व्याख्या
तशाच राहतात; आता आपण नवीन स्वयंसिद्धकेही मान्य करतो.

\begin{defn}[!!{identity} साठीची स्वयंसिद्धके]
\begin{align}
\ollabel{ax:id1} & \eq[t][t], \\
\ollabel{ax:id2} & \eq[t_1][t_2] \lif (!B(t_1) \lif
  !B(t_2)),
\end{align}
येथे $t$, $t_1$, $t_2$ ही कोणतीही बंद पदे आहेत.
\end{defn}

\begin{prop}\ollabel{prop:sound}
  \olref{ax:id1} आणि \olref{ax:id2} ही स्वयंसिद्धके वैध आहेत.
\end{prop}

\begin{proof}
  सराव.
\end{proof}

\begin{prob}
\olref[fol][axd][ide]{prop:sound} सिद्ध करा.
\end{prob}

\begin{prop}\ollabel{prop:iden1}
 कोणतेही पद $t$ आणि कोणताही संच~$\Gamma$ यांसाठी
 $\Gamma \Proves \eq[t][t]$.
\end{prop}

\begin{prop}\ollabel{prop:iden2}
  जर $\Gamma \Proves !A(t_1)$ आणि $\Gamma \Proves
  \eq[t_1][t_2]$, तर $\Gamma \Proves !A(t_2)$.
\end{prop}

\begin{proof}
पुढील !!{formula}
\[
(\eq[t_1][t_2] \lif (!A(t_1) \lif !A(t_2)))
\]
हा~\olref{ax:id2} चा दृष्टांत आहे. निष्कर्ष~\MP चे दोन उपयोग केल्यावर
मिळतो.
\end{proof}

\end{document}
